% =============================================================================
% Pré-processamento das bases de dados da tese
% Arquivo para \input{} no documento principal. Tradução: preprocessing.tex
% Pacotes usados: booktabs, url (ou hyperref), babel[brazilian]
% Manter sincronizado com scripts/build_thesis_datasets.py.
% Última atualização: 20/11/2026 (SIH 2024)
% =============================================================================

\section{Pré-processamento}
\label{sec:preprocessamento}

Cada base foi pré-processada por um script único e reprodutível
(\texttt{scripts/build\_thesis\_datasets.py}), que lê os arquivos originais sem
modificá-los e grava, junto com a base final, um relatório das exclusões etapa
por etapa. As exclusões são sequenciais: cada registro é contado apenas na
primeira etapa em que é removido.

% -----------------------------------------------------------------------------
\subsection{Sistema de Informações Hospitalares (SIH/SUS)}
\label{subsec:pre-sih}

\paragraph{Fonte e obtenção.}
Foram usados os arquivos de AIH reduzida (RD) do SIH/SUS, disponibilizados pelo
DATASUS em
\url{ftp://ftp.datasus.gov.br/dissemin/publicos/SIHSUS/200801_/Dados}. Cada
arquivo corresponde a uma unidade da federação e a uma competência, isto é, ao
mês em que a autorização de internação hospitalar (AIH) foi processada para
pagamento. Foram baixados os arquivos das 27 unidades da federação para as
competências de janeiro de 2024 a junho de 2025 (486 arquivos), em 19/11/2026,
pelo script \texttt{scripts/download\_datasus\_ftp.py}. Para cada arquivo, um
manifesto registra o caminho de origem, a data de extração, a data de
modificação no servidor, o tamanho e o \textit{hash} SHA-256, o que permite
verificar se uma reexecução usa exatamente os mesmos dados.

Os arquivos, em formato DBC, foram descompactados para DBF
(biblioteca \texttt{pyreaddbc}), e apenas 15 dos 113 campos da AIH foram lidos:
número da AIH (\texttt{N\_AIH}), tipo de AIH (\texttt{IDENT}), data de internação
(\texttt{DT\_INTER}), sexo (\texttt{SEXO}), raça/cor (\texttt{RACA\_COR}), idade e
unidade da idade (\texttt{IDADE}, \texttt{COD\_IDADE}), diagnóstico principal
(\texttt{DIAG\_PRINC}), caráter da internação (\texttt{CAR\_INT}), complexidade
(\texttt{COMPLEX}), especialidade do leito (\texttt{ESPEC}), municípios de
residência e do estabelecimento (\texttt{MUNIC\_RES}, \texttt{MUNIC\_MOV}), motivo
de saída (\texttt{COBRANCA}) e indicador de óbito (\texttt{MORTE}).

\paragraph{Unidade de análise e recorte temporal.}
A unidade de análise é a internação, identificada pela AIH. Foram selecionadas
as internações com data de internação em 2024. Como a competência indica o mês
de processamento e não o de internação, as AIHs de internações de 2024
continuam chegando nas competências de 2025: as competências de 2024 trouxeram
13.335.211 dessas AIHs e as de 2025 trouxeram outras 796.655, das quais 510.221
em janeiro e apenas 1.118 em junho. A leitura foi encerrada na competência de
junho de 2025, quando a chegada de novas AIHs referentes a 2024 já era
residual (menos de 0,01\% do total). Nenhuma internação de 2025 faz parte da
base: as competências de 2025 foram usadas apenas para completar o ano de 2024.

\paragraph{Critérios de exclusão.}
A Tabela~\ref{tab:sih-funil} resume as exclusões, aplicadas nesta ordem:
\begin{enumerate}
  \item \textbf{AIH reapresentada.} Uma mesma AIH pode ser processada em mais de
    uma competência (por exemplo, após rejeição e reapresentação). Manteve-se
    apenas o registro da competência mais recente de cada \texttt{N\_AIH}.
  \item \textbf{AIH de continuação de longa permanência} (\texttt{IDENT} = 5).
    Essas AIHs dão continuidade a uma internação já registrada em uma AIH
    principal e a contariam mais de uma vez.
  \item \textbf{Saída sem desfecho definido}, pelo primeiro dígito do motivo de
    saída (\texttt{COBRANCA}): permanência (2), em que o paciente continua
    internado; transferência para outro estabelecimento (3), em que o desfecho
    ocorre em outra internação; e encerramento administrativo (5).
  \item \textbf{Dados demográficos ou diagnóstico inválidos}: raça/cor sem
    informação (código 99), sexo diferente de masculino ou feminino, idade
    ausente ou acima de 120 anos e diagnóstico principal fora dos capítulos da
    CID-10. Apenas 131 internações de 2024 (0,001\%) foram excluídas nesta
    etapa, todas por idade acima de 120 anos. Os
    registros sem raça/cor observados nos arquivos são todos de AIHs de longa
    permanência de internações iniciadas antes de 2024.
\end{enumerate}

\begin{table}[htbp]
  \centering
  \caption{Exclusões aplicadas às AIHs com data de internação em 2024 (SIH/SUS).}
  \label{tab:sih-funil}
  \begin{tabular}{lrr}
    \toprule
    Etapa & Internações & \% do total \\
    \midrule
    AIHs com internação em 2024 & 14.131.866 & 100,00 \\
    AIH reapresentada em outra competência & $-$96.095 & 0,68 \\
    AIH de continuação de longa permanência & $-$49.822 & 0,35 \\
    Saída por permanência & $-$736.740 & 5,21 \\
    Saída por transferência & $-$630.287 & 4,46 \\
    Saída por encerramento administrativo & $-$155.443 & 1,10 \\
    Idade acima de 120 anos & $-$131 & 0,00 \\
    Raça/cor, sexo ou CID inválidos & 0 & 0,00 \\
    \midrule
    Internações mantidas & 12.463.348 & 88,19 \\
    \bottomrule
  \end{tabular}
\end{table}

\paragraph{Desfecho.}
O desfecho é binário: \emph{óbito} quando o indicador \texttt{MORTE} é igual a 1,
o que abrange as saídas por óbito (motivo 4) e os óbitos obstétricos (motivo 6),
e \emph{alta} nos demais casos, que, após as exclusões, correspondem às saídas
por alta (motivo 1) e às altas obstétricas (motivo 6). A taxa de óbito na base é
de 4,98\%.
Por ser derivado do registro de saída, o desfecho refere-se ao óbito durante a
internação, não ao óbito após a alta.

\paragraph{Atributos protegidos.}
Foram usados como atributos protegidos a raça/cor (Branca, Preta, Parda,
Amarela e Indígena), o sexo (Masculino e Feminino) e a faixa etária (0--17,
18--59 e 60 anos ou mais). A idade em anos foi obtida a partir de
\texttt{COD\_IDADE}: idades registradas em dias foram convertidas para
0 ano, em meses para o número de anos completos, em anos foram mantidas e o
código 5 (centenas de anos) corresponde a $100 + \texttt{IDADE}$. No SIH, a
raça/cor é registrada pelo estabelecimento de saúde e não necessariamente
autodeclarada pelo paciente, o que deve ser considerado na interpretação das
métricas de equidade.

\paragraph{Demais variáveis.}
Além dos atributos protegidos e do desfecho, a base guarda variáveis
conhecidas no momento da admissão: idade em anos, unidade da federação do
estabelecimento, categoria da CID-10 do diagnóstico principal (três
caracteres), caráter da internação (eletiva, urgência, acidente de trabalho,
acidente de trânsito ou outras lesões), complexidade (média ou alta),
especialidade do leito (cirúrgico, obstétrico, clínico, crônico, psiquiatria,
pneumologia sanitária, pediátrico, reabilitação ou hospital-dia/outros),
internação fora do município de residência (quando os códigos IBGE de seis
dígitos de \texttt{MUNIC\_RES} e \texttt{MUNIC\_MOV} diferem) e mês de
internação. Variáveis registradas durante ou após a internação, como uso de
UTI, dias de permanência, procedimentos realizados e valores pagos, foram
deliberadamente excluídas para evitar vazamento de informação sobre o
desfecho. O diagnóstico principal, embora usado como variável de entrada, é
consolidado no fechamento da AIH, o que constitui uma limitação.

\paragraph{Organização por capítulo da CID-10.}
Como o risco de óbito e o perfil dos pacientes variam muito entre grupos de
diagnóstico, as internações foram separadas pelos 22 capítulos da CID-10,
definidos pelas faixas oficiais de categorias de três caracteres, e cada
capítulo é analisado separadamente (Tabela~\ref{tab:sih-capitulos}). Foram
excluídos da análise os capítulos com taxa de óbito inferior a 0,2\%, cujo
desfecho não tem variação suficiente para a análise de viés, sobretudo nos
grupos interseccionais menores: XV (gravidez, parto e puerpério, 0,04\%), VII
(olho e anexos, 0,02\%) e VIII (ouvido e apófise mastoide, 0,20\%\footnote{Taxa
de 0,197\%, arredondada na tabela.}). Também foram excluídos os capítulos com
menos de 1.000 internações, que não permitem estimar métricas por grupo
interseccional: XXII (códigos para propósitos especiais, 415 internações). O
capítulo XX (causas externas, V01--Y98) não tem internações, pois seus códigos
não são aceitos como diagnóstico principal no SIH. Restam 17 capítulos. Os
números dos capítulos excluídos continuam registrados no relatório de
exclusões e no resumo da base.

\begin{table}[htbp]
  \centering
  \caption{Internações de 2024 por capítulo da CID-10 do diagnóstico principal (SIH/SUS).}
  \label{tab:sih-capitulos}
  \small
  \begin{tabular}{llrrr}
    \toprule
    Cap. & Descrição (categorias) & AIHs em 2024 & Mantidas & Óbito (\%) \\
    \midrule
    I     & Infecciosas e parasitárias (A00--B99)               &   965.970 &   842.452 & 14,95 \\
    II    & Neoplasias (C00--D48)                               & 1.104.508 & 1.002.953 &  8,02 \\
    III   & Sangue e transtornos imunitários (D50--D89)         &   138.202 &   120.828 &  5,51 \\
    IV    & Endócrinas, nutricionais e metabólicas (E00--E90)   &   257.186 &   215.324 &  5,95 \\
    V     & Transtornos mentais e comportamentais (F00--F99)    &   339.518 &   183.324 &  0,67 \\
    VI    & Sistema nervoso (G00--G99)                          &   254.320 &   208.990 &  4,64 \\
    VII$^{a}$ & Olho e anexos (H00--H59)                        &   198.289 &   194.476 &  0,02 \\
    VIII$^{a}$ & Ouvido e apófise mastoide (H60--H95)                &    30.450 &    28.954 &  0,20 \\
    IX    & Aparelho circulatório (I00--I99)                    & 1.332.909 & 1.081.637 &  9,93 \\
    X     & Aparelho respiratório (J00--J99)                    & 1.364.669 & 1.220.759 &  9,91 \\
    XI    & Aparelho digestivo (K00--K93)                       & 1.500.519 & 1.397.963 &  3,12 \\
    XII   & Pele e tecido subcutâneo (L00--L99)                 &   353.879 &   308.074 &  1,84 \\
    XIII  & Osteomuscular e tecido conjuntivo (M00--M99)        &   294.914 &   271.914 &  0,61 \\
    XIV   & Aparelho geniturinário (N00--N99)                   & 1.073.173 &   990.313 &  3,47 \\
    XV$^{a}$ & Gravidez, parto e puerpério (O00--O99)           & 2.106.315 & 1.999.809 &  0,04 \\
    XVI   & Afecções do período perinatal (P00--P96)            &   353.570 &   320.668 &  2,93 \\
    XVII  & Malformações congênitas (Q00--Q99)                  &    97.479 &    85.709 &  2,54 \\
    XVIII & Sintomas, sinais e achados anormais (R00--R99)      &   283.376 &   235.614 &  9,79 \\
    XIX   & Lesões e envenenamentos (S00--T98)                  & 1.582.739 & 1.291.275 &  2,56 \\
    XX$^{b}$ & Causas externas (V01--Y98)                       &         0 &         0 &   --- \\
    XXI   & Contato com serviços de saúde (Z00--Z99)            &   499.406 &   461.897 &  0,39 \\
    XXII$^{c}$ & Códigos para propósitos especiais (U00--U99)   &       475 &       415 &  7,23 \\
    \midrule
          & Total                                               & 14.131.866 & 12.463.348 & 4,98 \\
          & Total dos 17 capítulos analisados                   & 11.796.337 & 10.239.694 & 6,05 \\
    \bottomrule
    \multicolumn{5}{l}{\footnotesize $^{a}$ Excluído da análise: taxa de óbito inferior a 0,2\%.} \\
    \multicolumn{5}{l}{\footnotesize $^{b}$ Sem internações: códigos não aceitos como diagnóstico principal no SIH.} \\
    \multicolumn{5}{l}{\footnotesize $^{c}$ Excluído da análise: menos de 1.000 internações.} \\
  \end{tabular}
\end{table}

\paragraph{Arquivos gerados.}
Para cada capítulo analisado é gravado um arquivo Parquet
(\texttt{data/processed/thesis/sih\_v2/2024/sih\_v2\_2024\_capNN.parquet}), com as
variáveis categóricas codificadas como categorias; o conjunto ocupa cerca de
29~MB. Acompanham a base o relatório de exclusões por capítulo
(\texttt{sih\_v2\_2024\_exclusoes.csv}), a contagem de AIHs de 2024 por
competência (\texttt{sih\_v2\_2024\_cobertura\_competencias.csv}) e um resumo por
capítulo (\texttt{sih\_v2\_2024\_resumo.parquet}). A base pode ser regenerada com
\texttt{python scripts/download\_datasus\_ftp.py --inicio 2024-01 --fim 2025-06}
seguido de \texttt{python scripts/build\_thesis\_datasets.py --only sih}.
